\icode[magenta]{N22II-MM-CA-L4-260625-001}
\nceeID[20]{2022全国乙}{25}
如图（\(a\)），一质量为\(m\) 的物块\(A\)与轻质弹簧连接，静止在光滑水平面上；物块\(B\)向\(A\)运动，\(t = 0\)时与弹簧接触，到\(t = 2t_0\)时与弹簧分离，第一次碰撞结束，\(A\)、\(B\)的\(v - t\)图像如图（\(b\)）所示。已知从\(t = 0\) 到\(t = t_0\)时间内，物块\(A\)运动的距离为\(0.36v_0t_0\)。\(A\)、\(B\)分离后，\(A\)滑上粗糙斜面，然后滑下，与一直在水平面上运动的\(B\)再次碰撞，之后\(A\)再次滑上斜面，达到的最高点与前一次相同。斜面倾角为\(\theta\) （\(\sin\theta=0.6\)），与水平面光滑连接。碰撞过程中弹簧始终处于弹性限度内。求：
\begin{enumerate}
	\item 第一次碰撞过程中，弹簧弹性势能的最大值；
	\item 第一次碰撞过程中，弹簧压缩量的最大值；
	\item 物块\(A\)与斜面的摩擦因数。
\end{enumerate}

\begin{tikzpicture}[font=\scriptsize]
	\begin{scope}[rotate=37]
		\foreach \x in {0,.2,...,2.8}
			\draw[gray] (\x,0)--++(.2,-.2);
		\draw[thick](0,0)--(3,0);
	\end{scope}
	\foreach \x in {-4,-3.8,...,0}
		\draw[gray] (\x,0)--++(.2,-.2);
	\draw[thick] (-4,0)--(0,0);
	\draw[fill=orange!50] (-3.6,0) rectangle(-3.2,.4)node[midway]{\(B\)} (-1.5,0) rectangle(-1.1,.4)node[midway]{\(A\)};
	\draw[decoration={aspect=.6,segment length=3,amplitude=.9mm,coil},decorate] (-1.5,.2)--++(-.8,0);
	\node at(-.8,-.8){图（\(a\)）};
\end{tikzpicture}
\begin{tikzpicture}[font=\scriptsize,>=latex]
	\draw[->,blue] (-.5,0)--(5,0)node[below]{\(t /t_0\)};
	\draw[->,blue] (0,-.5)--(0,4)node[left]{\(v /v_0\)};
	\node at(-.2,-.2)[blue]{\(0\)};
	\draw[dashed,cyan] (0,3)node[left,blue]{ 2}--(4,3)--(4,0)node[below,blue]{ 2} (0,1.2)node[left,blue]{ 0.8}--(4,1.2) (0,1.5)node[left,blue]{ 1}--(2,1.5)--(2,0)node[below,blue]{ 1} (2,3)--(2,1.5);
	\draw[magenta] (0,0)..controls(.6,.05)and(1.5,.6)..(2,1.5)..controls(2.5,2.4)and(3.4,2.95)..(4,3)--++(.5,0)node[right]{\(A\)};
	\draw[teal] (0,1.8)node[left,blue]{\(1.2\)}..controls(.6,1.78)and(1.5,1.6)..(2,1.5)..controls(2.5,1.4)and(3.4,1.22)..(4,1.2)--++(.5,0)node[right]{\(B\)};
	\node at(2.25,-.8){图（\(b\)）};
\end{tikzpicture}
\begin{answer}
	（1）\(0.6mv_0^2\)（2）\(0.768v_0t_0\)（3）0.45
\end{answer}